\section{The arithmetic meaning of the geometry}
\label{sec:ipra}

\subsection{The base theory and the repository interface}

We use intuitionistic primitive recursive arithmetic with symbols generated by
zero, successor, projections, composition and primitive recursion, defining
equations for those symbols, equality axioms, successor separation and
injectivity, and induction for quantifier-free formulas.  In ProveIt's
\code{ConciseIPRA} development these ingredients are represented by the types
\code{PR k}, \code{PTm n}, \code{Fm}, the proof relation \code{Prf}, and the
theory \code{PRAAx}.  The derived rule \code{PRA.qfInd} requires an explicit
proof that its induction formula is quantifier free.  The files
\code{BiInterpretation.lean} and \code{NativeBiInterpretation.lean} provide
transport only after an actual object-language derivation has been built; an
ambient theorem over Lean's \code{Real} type does not automatically cross this
boundary.

All numerical data below are ultimately natural-number tuples.  Signed
integers may be coded as pairs \((p,n)\), denoting \(p-n\).  A rational is a
signed numerator together with a positive denominator.  Unreduced fractions
are sufficient: equality and order are tested by cross multiplication, and
all loops used below are bounded.

\subsection{Rapid dyadic names}

\begin{definition}[Regular rapid name]
A rapid dyadic name is a fixed primitive-recursive function
\(x:\N\to\Q\), together with a proof in \(\ipra\) that
\begin{equation}
  \abs{x(n+1)-x(n)}\le 2^{-n-4}                         \tag{2.1}
\end{equation}
for every \(n\), and an integer \(B_x\ge1\) satisfying
\(\abs{x(n)}\le B_x\).  The associated interval is
\begin{equation}
 I_x(n)=\bigl[x(n)-2^{-n-3},\,x(n)+2^{-n-3}\bigr].       \tag{2.2}
\end{equation}
\end{definition}

A finite geometric-sum induction gives, for \(m\ge n\),
\[
 \abs{x(m)-x(n)}<2^{-n-3}.
\]
Thus every later value lies in \(I_x(n)\).  Equality of names is the numerical
formula
\begin{equation}
 x=_\R y\quad:\Longleftrightarrow\quad
 \forall n\;\abs{x(n+6)-y(n+6)}\le2^{-n}.               \tag{2.3}
\end{equation}
The constants in (2.1)--(2.3) are deliberately generous.  Their purpose is to
make every precision shift primitive recursive and easy to audit.

\begin{definition}[Apartness and strict order]
An apartness witness \(x\apart0\) consists of natural numbers \(n,k\) and a
sign \(\sigma\in\{-1,+1\}\) such that the entire interval \(I_x(n)\) lies in
\([2^{-k},\infty)\) when \(\sigma=+1\), or in
\(({-}\infty,-2^{-k}]\) when \(\sigma=-1\).  A witness \(x<y\) is a finite
precision at which the upper endpoint of \(I_x\) is strictly below the lower
endpoint of \(I_y\).
\end{definition}

Apartness is stronger than \(x\ne0\) and is exactly the information needed to
compute a reciprocal.  This distinction is essential.  No proof below infers
a denominator bound from a bare negated equality.

\subsection{Arithmetic closure and bounded roots}

Addition and negation use fixed precision shifts.  On bounded inputs,
multiplication follows from
\begin{equation}
 \abs{xy-x'y'}\le \abs{x}\abs{y-y'}+\abs{y'}\abs{x-x'}.
                                                               \tag{2.4}
\end{equation}
If \(\abs{x}\ge d>0\) is supplied, reciprocal stability is
\begin{equation}
 \abs{x^{-1}-y^{-1}}
 =\frac{\abs{x-y}}{\abs{x}\abs{y}}
 \le \frac{2\abs{x-y}}{d^2}                                 \tag{2.5}
\end{equation}
whenever the input error is at most \(d/2\).  The right-hand side determines a
primitive-recursive precision transformer.

For a nonnegative rational \(r=p/q\), requested precision \(k\), and
\(S=2^k\), scan the bounded interval
\(0\le j\le(p+1)S\) and retain the last \(j\) with
\begin{equation}
 qj^2\le pS^2.                                                \tag{2.6}
\end{equation}
The endpoint fails (2.6), and the quantifier-free loop invariant says that the
retained value succeeds while its successor fails.  Hence
\[
 \frac jS\le\sqrt{p/q}<\frac{j+1}{S}.
\]
This is the square-root routine used for vector norms.  For nonnegative real
names, apply it to a sufficiently accurate rational enclosure.  The estimate
\begin{equation}
 \abs{\sqrt{x}-\sqrt{y}}\le\sqrt{\abs{x-y}}                 \tag{2.7}
\end{equation}
shows that doubling the requested input precision suffices near zero.  No
unbounded search and no exact real comparison occurs.

All finite loops just described have decidable rational guards.  Their
invariants can be expressed by primitive-recursive characteristic equations,
so quantifier-free induction applies.  Universal quantification over the final
precision parameter is introduced only after the finite estimate has been
proved.

\subsection{Constructive trigonometric names}
\label{subsec:trig-names}

The ProveIt trigonometric-identities packet supplies an analytic realization
that can be reused verbatim.  For a Gaussian rational \(z\) with
\(\norm{z}_1\le B\), use the finite exponential sum
\[
 E_N(z)=\sum_{j=0}^{N-1}\frac{z^j}{j!},
 \qquad N(B,n)=2B^2+2B+n+4.                                 \tag{2.8}
\]
Beyond index \(2B\), the majorant terms have ratio at most \(1/2\), and a
finite geometric estimate gives a tail below \(2^{-n-3}\).  A finite
Cauchy-product identity proves
\(\exp(z+w)=\exp z\exp w\).  Then
\begin{equation}
 \sin x=\frac{\exp(ix)-\exp(-ix)}{2i},\qquad
 \cos x=\frac{\exp(ix)+\exp(-ix)}2                         \tag{2.9}
\end{equation}
are primitive-recursive rapid names and satisfy the addition formulas and
\(\sin^2x+\cos^2x=1\).

A quarter turn \(q\in[1,2]\) is constructed by a bounded rational bisection
using explicit alternating-series signs for \(\cos1\) and \(\cos2\) and a
quantitative monotonicity estimate.  It satisfies
\begin{equation}
 \cos q=0,\qquad \sin q=1,
 \qquad \pi:=2q.                                             \tag{2.10}
\end{equation}
Thus every occurrence of \(\pi\), periodicity, or a supplementary angle in this
report has a concrete arithmetic realization.

\subsection{Angles from circle-valued data}

Most proofs need only a circle-valued angle pair \((c,s)\) satisfying
\(c^2+s^2=1\), with a sign witness for \(s\).  When an actual angle name is
desired, use two rationally separated charts.  Given a unit pair with
\(s>0\), compute once a rational approximation \(z\) to \(c\) with error at
most \(1/4\).
\begin{itemize}
\item If \(z\ge0\), then \(1+c\ge1/2\) and define
\[
 \theta=2\arctan\!\left(\frac{s}{1+c}\right).
\]
\item If \(z<0\), then \(1-c\ge1/2\) and define
\[
 \theta=\pi-2\arctan\!\left(\frac{s}{1-c}\right).
\]
\end{itemize}
The chart decision is finite and retained at every output precision.  The
arctangent is evaluated by range reduction and an alternating rational series
with an explicit primitive-recursive tail bound.  The Cayley identities
\[
 \cos(2\arctan t)=\frac{1-t^2}{1+t^2},\qquad
 \sin(2\arctan t)=\frac{2t}{1+t^2}                           \tag{2.11}
\]
prove that the resulting name realizes \((c,s)\).  For side lengths and
interior angles in \((0,\pi)\), the positive sine witness fixes the branch
uniquely.  This is a constructive \(\atanTwo(s,c)\) operation, not a global
continuous choice on the whole circle.

\subsection{The schema, not a universal evaluator}

For a fixed spherical triangle, the input consists of finitely many particular
\code{PR} syntax trees for coordinate names, bounds and certificates.  The
output operations are new particular \code{PR} trees obtained by composition,
primitive recursion and bounded searches.  The metatheory then constructs a
finite derivation
\[
 \Prf\;\PRAAx\;[]\;(\text{hypotheses}\to\text{conclusion}). \tag{2.12}
\]
There is no quantification over an arbitrary program code followed by an
unbounded universal evaluation operation.  Such an evaluator cannot enumerate
all primitive-recursive functions.  This fixed-program schema is the precise
sense in which the theorems are proved in first-order \(\ipra\).

\begin{remark}[Ambient Lean versus object-level arithmetic]
A proof by \code{ring} over a Lean commutative ring is valuable: it checks the
finite identity that the arithmetic derivation must replay.  It is not itself
a term of \(\Prf\;\PRAAx\).  The final step must either write the corresponding
equational derivation directly or invoke a polynomial normalizer whose
soundness has itself been proved inside the object theory.  The Lean roadmap in
Section~\ref{sec:formalization} preserves this distinction explicitly.
\end{remark}
